\chapter{中值定理、证明题与不等式}
\section{定理条件与使用信号}
\ruleitem{费马定理}{若 $x_0$ 是定义区间内部的局部极值点且 $f'(x_0)$ 存在，则 $f'(x_0)=0$。驻点未必是极值，$x^3$ 在 0 即为反例；端点极值不要求导数为零。}
\ruleitem{Rolle}{若 $f\in C[a,b]$、在 $(a,b)$ 可导，且 $f(a)=f(b)$，则存在 $\xi\in(a,b)$ 使 $f'(\xi)=0$。有 $m+1$ 个不同零点，反复使用可得到相应导数的零点。}
\ruleitem{Lagrange}{同样的连续、可导条件下，存在 $\xi\in(a,b)$：$f(b)-f(a)=f'(\xi)(b-a)$。用来把“函数差”变成“导数界”。}
\ruleitem{Cauchy}{若 $f,g$ 均满足上述条件，存在 $\xi$：
$[f(b)-f(a)]g'(\xi)=[g(b)-g(a)]f'(\xi)$。
要写比值型还需相应分母非零，常用 $g'(x)\ne0$ 保证。}
\ruleitem{L'Hôpital}{当分子、分母同趋于 0，或分母绝对值趋于无穷的相应型，函数在去心单侧邻域可导、$g'\ne0$，且 $f'/g'$ 有有限或无穷极限时，可推出 $f/g$ 的该极限。$0\cdot\infty$、$\infty-\infty$、幂指型须先改写。导数比极限不存在，不代表原比值不存在。}

\textbf{证明题的四个入口。}
目标有 $f'$ 的零点，考虑 Rolle；有 $f'+pf$，考虑乘积分因子 $\ee^{\int p}$；有两个不同点的导数，先找中间分点再在两侧分别用中值定理；有积分与函数值，先把积分变为原函数。

\section{单调、凹凸、零点与基本不等式}
$f'\ge0$ 可推出不减，$f'>0$ 推出严格增；反向不成立，例如 $x^3$ 严格增却在 0 导数为 0。二阶导数 $f''\ge0$ 时称为凸函数，本文采用弦在图像上方的约定：
\[
f(\lambda x+(1-\lambda)y)\le\lambda f(x)+(1-\lambda)f(y),\quad0\le\lambda\le1.
\]
可微凸函数位于每条切线之上；凹函数不等号反向。拐点要求凹凸性改变，$f''=0$ 仅是候选条件。闭区间求最值需比较驻点、不可导点及端点。
\begin{align*}
\ee^x&\ge1+x &&(x\in\R),\\
\ln(1+x)&\le x &&(x>-1),\\
|\sin x|&\le |x| &&(x\in\R),\\
2|ab|&\le \lambda a^2+\lambda^{-1}b^2 &&(\lambda>0),\\
\left|\sum a_kb_k\right|^2&\le\left(\sum a_k^2\right)\left(\sum b_k^2\right),\\
\left|\int_a^b uv\dd x\right|^2&\le\int_a^b u^2\dd x\int_a^b v^2\dd x.
\end{align*}
Cauchy 不等式的等号要求两列成比例；积分版本对连续函数要求处处成比例（一般可积情形为几乎处处）。Jensen 不等式是凸性向有限加权和的推广；使用前检查权重非负、和为 1 及凸性区间。

\ruleitem{根的数量}{存在用介值；唯一用严格单调；最多 $m$ 个常用反证：若有 $m+1$ 个不同根，经 $m$ 次 Rolle 迫使 $f^{(m)}$ 有零点。约束导数的符号有助于排除多根。}
\ruleitem{高阶导数判别}{若 $f'(a)=\cdots=f^{(m-1)}(a)=0$，$f^{(m)}(a)\ne0$，则首个非零阶 $m$ 为偶数时按该导数正负判断极小、极大；奇数时不是极值。需有相应 Taylor 展开。}

\section{零点、辅助函数与积分不等式真题}
\problem{22}{严格凸函数恰有两个零点}{2010 年第 2 届，非数学类，二，\src{02}}
设 $f''(x)>0$，$\lim_{x\to+\infty}f'(x)=\alpha>0$，$\lim_{x\to-\infty}f'(x)=\beta<0$，且某点 $f(x_0)<0$。证明恰有两个实根。
\sol 充分大的 $x$ 有 $f'(x)>\alpha/2$，积分或中值定理给出 $f(x)\to+\infty$。充分小的 $x$ 有 $f'(x)<\beta/2<0$；向负方向积分同样得到 $f(x)\to+\infty$。与 $f(x_0)<0$ 配合，左右两侧各有一根。
若有三个不同根，两次 Rolle 使 $f''$ 在某点为 0，与严格正矛盾，故恰有两根。
\trap{负半轴的导数为负，向左移动时函数反而升高。不要把 $x\to-\infty$ 时的符号判反。}

\problem{23}{端点条件迫使三阶导数取指定值}{2011 年第 3 届，非数学类，三，\src{03}}
$f\in C^3[-1,1]$，$f(-1)=0$，$f(1)=1$，$f'(0)=0$。证明存在 $\xi\in(-1,1)$，$f'''(\xi)=3$。
\sol 在 0 作二阶 Taylor 展开，余项取三阶：
\[1=f(0)+\tfrac12f''(0)+\tfrac16f'''(\eta_+),\qquad
0=f(0)+\tfrac12f''(0)-\tfrac16f'''(\eta_-),\]
其中 $\eta_-\in(-1,0)$，$\eta_+\in(0,1)$。两式相减得 $f'''(\eta_+)+f'''(\eta_-)=6$。连续函数 $f'''$ 在两值之间必取平均数 3（若两值已等于 3，直接取该点）。
\trap{两次 Taylor 余项的中间点一般不同，不能擅自合并为同一个点。这里靠连续介值完成最后一步。}

\problem{24}{两个不同点的倒数关系}{2016 年第 8 届，非数学类，五，\src{08}}
连续函数 $f$ 满足 $I=\int_0^1 f\dd x\ne0$。证明存在不同的 $x_1,x_2\in(0,1)$，使 $1/f(x_1)+1/f(x_2)=2/I$。
\method{先找“累计积分的一半”，再在两段分别用中值定理。}
\sol 令 $F(x)=I^{-1}\int_0^x f(t)\dd t$，$F(0)=0,F(1)=1$。介值定理给出 $c\in(0,1)$、$F(c)=1/2$。在 $[0,c]$ 与 $[c,1]$ 分别用 Lagrange，得到
\[f(x_1)=\frac I{2c},\quad f(x_2)=\frac I{2(1-c)},\quad 0<x_1<c<x_2<1.\]
两值均非零，倒数相加即为 $2/I$。$f$ 不必恒正，$F$ 不必单调；介值已经足够。

\problem{25}{先介值，再两次中值}{2020 年第 12 届，非数学类，三，\src{12}}
$f\in C[0,1]$，在 $(0,1)$ 可导，$f(0)=0,f(1)=1$。证明存在 $c\in(0,1)$ 使 $f(c)=2-3c$；并存在不同 $\xi,\eta\in(0,1)$，使 $(1+f'(\xi))(1+f'(\eta))=4$。
\sol 对 $g(x)=f(x)+3x-2$，有 $g(0)=-2,g(1)=2$，故存在该 $c$。分别在两侧使用 Lagrange：
\[f'(\xi)=\frac{2-3c}{c},\qquad f'(\eta)=\frac{3c-1}{1-c},\quad \xi<c<\eta.\]
于是 $1+f'(\xi)=2(1-c)/c$，$1+f'(\eta)=2c/(1-c)$，乘积为 4。

\problem{26}{由零积分到非线性导数等式}{2023 年第 15 届，非数学 B 类，五，\src{15B}}
$f$ 在 $[0,1]$ 可导，$f(0)>0,f(1)>0,\int_0^1 f=0$。证明至少有两个零点，且某 $\xi\in(0,1)$ 满足 $f'(\xi)+3f^3(\xi)=0$。
\sol 若 $f\ge0$，由端点附近连续且为正，积分必严格大于零。因此有 $x_0\in(0,1)$ 使 $f(x_0)<0$。两侧各由介值定理得零点 $a<b$。
构造 $G(x)=f(x)\exp(\int_0^x3f^2(t)\dd t)$，则 $G(a)=G(b)=0$。Rolle 给出 $G'(\xi)=0$，而
\[G'(x)=\exp\left(\int_0^x3f^2(t)\dd t\right)[f'(x)+3f^3(x)],\]
指数因子非零，所证即得。看到 $f'+3f^3=f'+(3f^2)f$，就能识别积分因子。

\problem{27}{分部积分加 Cauchy，连等号也要证明}{2023 年第 15 届，非数学 A 类五／B 类六，共题，\src{15A}\src{15B}}
$f\in C^1[0,1]$，$f(0)=0$。证明 $\int_0^1 f^2\le4\int_0^1(1-x)^2(f')^2$，并求等号条件。
\sol 记两积分为 $I,J$。分部积分得
\[I=[(x-1)f^2]_0^1+2\int_0^1(1-x)ff'\dd x
=2\int_0^1(1-x)ff'\dd x\le2\sqrt{IJ}.\]
$I=0$ 时 $f\equiv0$；$I>0$ 时平方除以 $I$ 得 $I\le4J$。
若取等且 $I>0$，Cauchy 等号要求 $(1-x)f'=cf$。代回 $I=2\int(1-x)ff'$ 得 $c=1/2$。于是 $f=C(1-x)^{-1/2}$（$x<1$），由 $f(0)=0$ 得 $C=0$，与 $I>0$ 矛盾。所以等号当且仅当 $\ans{f\equiv0}$。

\problem{28}{导数满足 Lipschitz 条件的全局估计}{2024 年第 16 届，非数学 A 类五／B 类六，共题，\src{16A}\src{16B}}
$f\in C^1(\R)$，$f\ge0$，且 $|f'(x)-f'(y)|\le M|x-y|$，$M>0$。证明 $(f'(x))^2\le2Mf(x)$。
\sol 由微积分基本定理，对任意实数 $h$，
\[|f(x+h)-f(x)-f'(x)h|\le\frac M2h^2.\]
故 $0\le f(x+h)\le f(x)+f'(x)h+Mh^2/2$。右端二次式对每个 $h$ 均非负，取 $h=-f'(x)/M$，得 $0\le f(x)-(f'(x))^2/(2M)$。
\trap{定义域是整个实轴，所以可自由取最优 $h$。若只在有限区间非负，不能直接照搬此选择。}

\problem{29}{把目标不等式变成两个单调函数}{2016 年第 8 届，非数学类，二，\src{08}}
$f$ 在 $[0,1]$ 可导，$f(0)=0$，$0<f'(x)<1$。证明对 $0<a<1$，$(\int_0^a f)^2>\int_0^a f^3$。
\sol 对 $x>0$，$f(x)>0$。令 $A(x)=\int_0^x f(t)\dd t$，
$G(x)=2A(x)-f^2(x)$。因 $G(0)=0$、$G'=2f(1-f')>0$，有 $G(x)>0$。
再令 $H(x)=A^2(x)-\int_0^x f^3(t)\dd t$，则 $H(0)=0$，
\[H'(x)=2A(x)f(x)-f^3(x)=f(x)G(x)>0.\]
故 $H(a)>0$。目标里同时出现“积分的平方”和“三次方积分”，先整体求导，再对剩余括号求导，是可迁移的构造路径。
